\documentclass[11pt]{article}
\usepackage[margin=1in]{geometry}
\usepackage[T1]{fontenc}
\usepackage[utf8]{inputenc}
\usepackage{amsmath,amssymb,amsthm}
\usepackage{enumitem}

\title{ICPC World Finals 2025\\E. Delivery Service}
\author{}
\date{}

\begin{document}
\maketitle

\section*{Problem Summary}

After each courier is hired, we must count how many unordered pairs of cities $(u,v)$ can send packages
to each other. The time schedule is periodic, so the natural model is a two-layer graph: one copy of each
city for the ``home-city time'' and one copy for the ``destination-city time''.

\section*{Key Observations}

\begin{itemize}[leftmargin=*]
    \item Create two vertices for every city $u$:
    $u_L$ for the times when couriers are in their home city, and $u_R$ for the times when couriers are in
    their destination city.
    A courier $(a,b)$ creates an undirected edge between $a_L$ and $b_R$.
    \item A city is represented by the unordered pair of connected components containing its two copies.
    For example, if both copies lie in the same component $C$, the city corresponds to $(C,C)$.
    If the copies lie in two different components $C$ and $D$, the city corresponds to $(C,D)$.
    \item Two cities $u$ and $v$ are connected if and only if their unordered component pairs
    \emph{intersect}. Equivalently, the two cities share at least one DSU component among their four
    copies.
    \item Therefore, if for each component $C$ we know how many cities touch $C$, then
    $\binom{\text{touch}(C)}{2}$ counts all city pairs that share component $C$.
    A pair of cities can be counted twice only when both are split across the same two components
    $(C,D)$, so those double counts must be subtracted once.
    \item We can maintain this incrementally with a DSU.
    For each DSU component we store how many cities have both copies already inside that component, and
    a hash map counting how many cities are split between this component and another component.
    When two DSU components merge, only counters touching these two components change.
\end{itemize}

\section*{Algorithm}

\begin{enumerate}[leftmargin=*]
    \item Build a DSU on $2n$ vertices, one left copy and one right copy for each city.
    Initially, every city is split between its two singleton components, so we insert one count for the
    component pair $(u_L,u_R)$.
    \item Maintain two global totals:
    \begin{itemize}[leftmargin=*]
        \item \texttt{tot1}: the sum over all components $C$ of
        $\binom{z_C}{2}$, where $z_C$ is the number of cities that have at least one copy in $C$;
        \item \texttt{tot2}: the sum over all unordered component pairs $(C,D)$ with $C \ne D$ of
        $\binom{y_{C,D}}{2}$, where $y_{C,D}$ is the number of cities split between $C$ and $D$.
    \end{itemize}
    The required answer is \texttt{tot1 - tot2}.
    \item When a courier $(a,b)$ is added, unite the DSU components containing $a_L$ and $b_R$.
    Merge the smaller hash map into the larger one and update the two global totals in $O(\text{moved map
    entries})$ expected time.
\end{enumerate}

\section*{Correctness Proof}

We prove that the algorithm returns the correct answer.

\paragraph{Lemma 1.}
Two cities $u$ and $v$ are connected if and only if the unordered pair of DSU components containing
$(u_L,u_R)$ intersects the unordered pair containing $(v_L,v_R)$.

\paragraph{Proof.}
Suppose the two unordered pairs intersect, and let $C$ be a common component.
Then one copy of $u$ and one copy of $v$ both lie in $C$, so a package can be transferred through the
sequence of couriers represented by paths inside $C$.
Because the couriers are reversible over successive half-days, the same argument also gives a route in the
opposite direction. Hence the cities are connected.

Conversely, if the two unordered pairs are disjoint, then no copy of $u$ shares a connected component
with any copy of $v$. No sequence of handoffs can ever move a package from a copy of $u$ to a copy of
$v$, so the cities are not connected. \qed

\paragraph{Lemma 2.}
At every moment, the DSU data stored by the program represents exactly:
\begin{itemize}[leftmargin=*]
    \item for each component $C$, how many cities touch $C$;
    \item for each unordered pair $(C,D)$ with $C \ne D$, how many cities are split between $C$ and $D$.
\end{itemize}

\paragraph{Proof.}
Initially this is true by construction: each city touches exactly its two singleton components.
A DSU union only changes data that involve one of the two merged components.
The program removes the old contributions of those counts and inserts the new merged contributions.
All other cities keep exactly the same touched components and split-component pair. Therefore the
maintained counts are always exact. \qed

\paragraph{Theorem.}
The algorithm outputs the correct answer for every valid input.

\paragraph{Proof.}
By Lemma 2, after each courier is added the program knows the exact values of all $z_C$ and $y_{C,D}$.
Summing $\binom{z_C}{2}$ over all components counts every pair of cities once for each shared
component. By Lemma 1, a pair is connected exactly when it shares at least one component.
A pair can be counted twice only if both cities are split across the same two components $(C,D)$, and the
subtracted term $\binom{y_{C,D}}{2}$ removes exactly this double counting.
Therefore \texttt{tot1 - tot2} is exactly the number of connected city pairs. \qed

\section*{Complexity Analysis}

Each DSU operation is nearly constant aside from hash-map merging. With the usual small-to-large
strategy, every map entry is moved only $O(\log n)$ times, so the total expected running time is
$O((n+m)\log n)$ and the memory usage is $O(n+m)$.

\section*{Implementation Notes}

\begin{itemize}[leftmargin=*]
    \item The code stores only counts for split cities; cities whose two copies are already in the same DSU
    component are tracked by a separate counter per component.
    \item Hash maps are merged from the smaller component into the larger one to keep the total work low.
\end{itemize}

\end{document}
